\section{Domini numerici e interpretazione astratta}

\subsection{Il problema generale}

\begin{nota}{Problemi di affidabilità del software}{}
L'hardware è diventato milioni di volte più potente rispetto a 25 anni fa, e
la dimensione dei programmi è cresciuta in modo simile, raggiungendo decine
di milioni di righe di codice.
Questi programmi devono essere progettati, sviluppati e mantenuti per lunghi
periodi (fino a 20 anni e più) a costi ragionevoli.

I team di sviluppo e manutenzione non assistiti non possono far fronte a
questa esplosione di dimensioni e complessità. Molti software presentano un
numero di bug a volte difficilmente sopportabile anche nelle applicazioni da
ufficio.
Nessuna applicazione critica per la sicurezza può tollerare questo tasso di
fallimento.

Il problema dell'affidabilità del software è uno dei problemi più importanti
che l'informatica deve affrontare.
Questo giustifica il crescente interesse per strumenti meccanici che aiutino
il programmatore a ragionare sui programmi.
\end{nota}

\subsection{Esempi di problemi di verifica}

\begin{esempio}{Disastro di Ariane 5}{}
Il disastro del volo 501 di Ariane 5 è stato causato da un segmento di
programma che convertiva un numero in virgola mobile in un intero a 16 bit
con segno.
Il programma è stato eseguito con un valore di input al di fuori
dell'intervallo rappresentabile da un intero a 16 bit con segno.
Questo errore di runtime (overflow), si è verificato sia nei computer attivi
che in quelli di backup quasi contemporaneamente, ed entrambi i computer si
sono spenti.
Ciò ha comportato la perdita totale del controllo dell'assetto. L'Ariane 5
è diventato incontrollabile e le forze aerodinamiche hanno distrutto il
veicolo.
\end{esempio}

\begin{esempio}{Disastri storici}{}
\begin{itemize}
    \item Il Mars Climate Orbiter si è disintegrato nell'atmosfera marziana
    nel 1999 dopo aver mancato l'inserimento in orbita a causa di calcoli
    incoerenti sulle unità di misura.
    \item Nello stesso anno, si sospetta che il Mars Polar Lander si sia
    schiantato su Marte all'atterraggio perché un flag software non è stato
    resettato correttamente.
    \item Nella missione dimostrativa tecnologica Mars Pathfinder (MPF) del
    1997, un giorno di esplorazione è andato perso quando i team di supporto a
    terra sono stati costretti a riavviare il sistema durante il download dei
    dati scientifici.
    \item La missione Mars Exploration Rover (MER) della NASA del 2003
    comprende due rover. Con un costo di 400 milioni di dollari per ogni
    rover, un errore di codifica che spegne un rover durante la notte sarebbe
    a tutti gli effetti un errore da 4,4 milioni di dollari, oltre a una
    perdita di tempo prezioso di esplorazione sul pianeta.
\end{itemize}
\end{esempio}

\begin{esempio}{È ben definito $x/(x-y)$?}{}
Molte cose possono andare storte:
\begin{itemize}
    \item $x$ e/o $y$ potrebbero non essere inizializzati;
    \item $x-y$ potrebbe andare in overflow;
    \item $x$ e $y$ potrebbero essere uguali (o $x-y$ potrebbe andare in underflow):
    divisione per 0;
    \item $x/(x-y)$ potrebbe andare in overflow (o underflow).
\end{itemize}
\end{esempio}

\subsection{Metodi formali di verifica dei programmi}

\begin{definizione}{Scopo della verifica formale}{}
Lo scopo della verifica formale è dimostrare meccanicamente che tutte le
possibili esecuzioni di un programma sono corrette in tutti gli ambienti di
esecuzione specificati.
\end{definizione}

Esistono diversi metodi di verifica:
\begin{itemize}
    \item Metodi deduttivi.
    \item Model checking.
    \item Tipizzazione dei programmi.
    \item Analisi statica.
\end{itemize}

\begin{nota}{L'indecidibilità sui metodi di verifica}{}
A causa dell'indecidibilità della verifica dei programmi, tutti i metodi
sono parziali o incompleti e ricorrono a una qualche forma di
approssimazione.
\end{nota}

\subsection{Interpretazione astratta}

\begin{definizione}{Interpretazione astratta}{}
L'interpretazione astratta è il framework giusto per lavorare con il concetto
di approssimazione corretta.
È una teoria per approssimare insiemi e operazioni su insiemi come
considerato nella teoria degli insiemi (o delle categorie), comprese le
definizioni induttive.
È una teoria dell'approssimazione del comportamento dei sistemi discreti
dinamici.
La computazione avviene su un dominio di proprietà astratte: il dominio
astratto, utilizzando operazioni astratte che sono approssimazioni corrette
delle operazioni concrete.
La correttezza segue dal progetto.

\end{definizione}

\begin{nota}{Livelli di approssimazione}{}

L'astrazione (approssimazione) può essere abbastanza grossolana da essere
calcolabile finitamente, ma abbastanza precisa da essere praticamente utile.

\end{nota}

\subsection{Domini astratti numerici}

\begin{definizione}{Domini astratti numerici}{}
I domini astratti numerici sono utilizzati per rappresentare proprietà
numeriche delle variabili di un programma.
Alcuni esempi di domini astratti numerici sono:
\begin{itemize}
    \item Segni
    \item Bounding Boxes
    \item Congruenze semplici
    \item Differenze limitate
    \item Ottagoni
    \item Poliedri convessi
    \item Z-Polyhedra
    \item Congruenze relazionali
    \item Coni polinomiali
\end{itemize}
\end{definizione}

\subsection{Un assaggio di semantica concreta e astratta}

\begin{nota}{Semantica Concreta vs. Astratta}{}
La semantica concreta descrive l'esatto comportamento di un programma,
mentre la semantica astratta ne descrive un comportamento approssimato.
La semantica concreta opera su un dominio concreto (es. $\wp(\mathbb{R}^2)$),
mentre la semantica astratta opera su un dominio astratto (es. poliedri
convessi).
\end{nota}

\begin{esempio}{Semantica concreta}{}
Considerando il seguente programma:
\begin{verbatim}
x := 0; y := 0;
while x <= 100 do
    read(b);
    if b then x := x+2
    else x := x+1; y := y+1;
    endif
endwhile
\end{verbatim}
La semantica concreta calcola l'insieme esatto di tutti i possibili valori
di (x, y) in ogni punto del programma.
\end{esempio}

\begin{esempio}{Semantica astratta}{}
La semantica astratta calcola un'approssimazione dell'insieme di tutti i
possibili valori di (x, y) in ogni punto del programma.
Per garantire la terminazione, vengono utilizzati operatori di widening per
accelerare la convergenza del punto fisso.
\end{esempio}
